\chapter{\Localized{Methodology}{方法}}
\label{chap:methodology}

\Localized{This chapter demonstrates common scientific elements in the same visual system. English, Chinese, mathematics, tables, and code can coexist without local font commands.}{本章展示同一视觉体系下常见的科学内容。英文、中文、数学公式、表格和代码可以直接混排，无需在正文中手动切换字体。}

\section{\Localized{Model Configuration}{模型配置}}

\Localized{Table~\ref{tab:methodology:configuration} demonstrates the mandatory full-width table contract. All manuscript tables use the AcademicTable interface; column widths are distributed across the complete text block.}{表~\ref{tab:methodology:configuration} 展示强制满行宽表格规范。所有正文表格均通过 AcademicTable 接口生成，各列在完整正文宽度内分配。}

\begin{table}[tb]
  \centering
  \caption{\Localized{Example computational configuration.}{示例计算配置。}}
  \label{tab:methodology:configuration}
  \begin{AcademicTable}{LLL}
    \toprule
    \Localized{Component}{组件} & \Localized{Choice}{选择} & \Localized{Purpose}{用途} \\
    \midrule
    \Localized{Solver}{求解器} & ODE & \Localized{Dynamic integration}{动态积分} \\
    \Localized{Backend}{后端} & PyTorch & \Localized{Differentiable computation}{可微计算} \\
    \Localized{Device}{设备} & MPS & \Localized{Apple Silicon acceleration}{Apple Silicon 加速} \\
    \bottomrule
  \end{AcademicTable}
\end{table}

\section{\Localized{Scientific Code}{科学代码}}

\Localized{The code environment uses Maple Mono explicitly; compilation stops when the required font asset is unavailable.}{代码环境显式使用 Maple Mono；当所需字体文件不存在时，编译会直接停止。}

\begin{lstlisting}[language=Python,caption={Minimal process-rate function.},label={lst:methodology:rate}]
def process_rate(x, theta):
    """Return a Monod-type process rate."""
    mu_max, half_saturation = theta
    return mu_max * x / (half_saturation + x)
\end{lstlisting}

\section{\Localized{Minimal Figure}{最小示意图}}

\Localized{Figure~\ref{fig:methodology:workflow} is constructed directly in LaTeX so the template has no mandatory external figure dependency.}{图~\ref{fig:methodology:workflow} 直接由 LaTeX 构建，因此模板本身不依赖外部示意图文件。}

\begin{figure}[tb]
  \centering
  \setlength{\unitlength}{1mm}
  \begin{picture}(120,34)
    \put(2,10){\framebox(30,14){Scientific data}}
    \put(45,10){\framebox(30,14){Model}}
    \put(88,10){\framebox(30,14){Prediction}}
    \put(32,17){\vector(1,0){13}}
    \put(75,17){\vector(1,0){13}}
  \end{picture}
  \caption{\Localized{Minimal scientific modelling workflow used for layout demonstration.}{用于版式展示的最小科学建模流程。}}
  \label{fig:methodology:workflow}
\end{figure}
